\section{Artifact Appendix}
\label{app:artifact_alghard}

\subsection{Abstract}
The artifact contains several components. The algorithmic
component reproduces the Table~\ref{tab:alg} accuracy results for vanilla RSD,
backtracking-enhanced RSD (BRSD), and BRSD with optimizations. The hardware
component provides a trace-driven simulator and plotting workflow for
reproducing the Figure~\ref{fig:5_end2end} latency, goodput, and energy results from packaged
CSV/JSONL inputs without rerunning model inference.

\subsection{Artifact Check-list}
\begin{itemize}[leftmargin=*, itemsep=1pt]
  \item \textbf{Algorithm:} speculative reasoning with RSD, BRSD, and BRSD with optimized execution.
  \item \textbf{Program:} Python scripts, shell runners, step-profile CSVs, JSONL token logs, reference result summaries, and plotting utilities.
  \item \textbf{Models:} Qwen2.5 0.5B draft, Qwen2.5-Math-PRM-7B, Qwen2.5 7B target, and Qwen2.5 1.5B target. Model checkpoints are not redistributed.
  \item \textbf{Datasets:} Math500, GSM8K, Gaokao2023EN, and OlympiadBench.
  \item \textbf{Compute platform:} Algorithmic runs use 3$\times$ NVIDIA RTX 3090 GPUs in the reference setup. The packaged hardware simulator and Figure~\ref{fig:5_end2end} plotting workflow run from CSV/JSONL traces and do not require GPUs.
  \item \textbf{Expected time:} The algorithmic smoke test takes a few minutes after model initialization; the Config1/Math500 key result takes about 1--2 hours; the full algorithmic workflow can take about 1--2 days depending on GPU availability. The packaged hardware plotting workflow takes seconds, and regenerating simulator CSVs from the packaged traces takes minutes.
  
\end{itemize}

\subsection{Description}
The artifact codebase: \url{https://zenodo.org/records/22015981} or \url{https://github.com/zehuanzhang/HeteroReason}.
The artifact is organized into several directories. \texttt{AE\_alg}
contains the algorithmic reproduction workflow under \texttt{Algorithm}.
\texttt{AE\_hardware} contains the communication-aware trace-driven
latency/energy simulator, packaged step-profile CSVs and logs, and precomputed
simulator outputs. \texttt{AE\_hardware\_V80} is the V80 projection of the
same simulator. The root-level \texttt{figure8} directory contains the final
four-bar plotting script and its packaged metric inputs.

The three Table~\ref{tab:alg} algorithmic methods are:
\begin{itemize}[leftmargin=*, itemsep=1pt]
  \item \texttt{rsd}: the vanilla speculative reasoning baseline.
  \item \texttt{brsd}: RSD with backtracking.
  \item \texttt{brsd\_optimized}: BRSD with the optimizations.
\end{itemize}
These correspond to the implementation names \texttt{beam1},
\texttt{bbeam1}, and \texttt{bbeam1\_prefetch\_cache}, respectively.

\subsection{Installation}
After downloading the artifact, install the Python dependencies and set local
model paths for algorithmic reproduction:
{\footnotesize
\begin{verbatim}
pip install -r AE_alg/requirements.txt
export DRAFT_MODEL=/path/to/Qwen2.5-0.5B-Instruct
export PRM_MODEL=/path/to/Qwen2.5-Math-PRM-7B
export CONFIG1_TARGET_MODEL=/path/to/Qwen2.5-7B-Instruct
export CONFIG2_TARGET_MODEL=/path/to/Qwen2.5-1.5B-Instruct
export CUDA_VISIBLE_DEVICES=0,1,2
\end{verbatim}
}
The hardware simulator and Figure~\ref{fig:5_end2end} plotting workflow use packaged CSV/JSONL
inputs and do not require model checkpoints.

\subsection{Algorithmic Experiment Workflow}
The smoke test validates the environment by running the three Table~\ref{tab:alg} methods
on eight fixed Config1/Math500 examples:
\begin{verbatim}
cd AE_alg/Algorithm
bash run_smoke.sh
\end{verbatim}

The key-result workflow focuses on Config1/Math500. The default key-result
command runs \texttt{brsd\_optimized}; reviewers may also run the three methods
separately:
\begin{verbatim}
cd AE_alg/Algorithm
bash run_key_results.sh

KEY_METHODS="rsd" bash run_key_results.sh
KEY_METHODS="brsd" bash run_key_results.sh
KEY_METHODS="brsd_optimized" bash run_key_results.sh
\end{verbatim}

The full algorithmic workflow reproduces Table~\ref{tab:alg} across both configurations,
all four datasets, and all three methods:
\begin{verbatim}
cd AE_alg/Algorithm
bash run_full_results.sh
\end{verbatim}

\subsection{Hardware Simulator and Figure Workflow}
The hardware component can reproduce Figure~\ref{fig:5_end2end} without
rerunning model inference:
\begin{verbatim}
python figure8/plot_figure8.py --normalized
\end{verbatim}


\subsection{Evaluation and Expected Results}
The algorithmic smoke test reports per-method accuracy and latency on the fixed
eight-example subset. The key-result workflow reproduces the Config1/Math500
Table~\ref{tab:alg} entries, with \texttt{brsd\_optimized} as the recommended artifact
result. The full workflow reproduces all Table~\ref{tab:alg} algorithmic accuracy values
across Config1 and Config2.

The hardware workflow regenerates Figure~\ref{fig:5_end2end} and the corresponding
CSV/JSON metric summaries. The
expected metrics include average latency, goodput, and average energy per
problem.

\subsection{Notes}
Algorithmic latency is platform-dependent and can vary with GPU
types. The hardware simulator
and plotting workflow are deterministic for the supplied CSV/JSONL inputs. Model
checkpoints are referenced by local paths or Hugging Face identifiers and are
not redistributed with the artifact.
